\section{Proxy-Corrected Local-Law Objective}
\label{app:v8-objective-ipw}

The local laws are oracle statements: a produced state should preserve the
value assigned by \(f^\ast\). In the LLM setting, most local-law losses are
observed through a proxy \(\widehat f\), while direct \(f^\ast\) node labels are
available only on sampled nodes. This appendix spells out the adjustment used
in Section~\ref{sec:v8-objective}.

\subsection{Local-Law Estimand}
\label{app:v8-local-law-estimand}

For a realized tree \(T\) with node set \(V(T)\), let \(g\) be the single
learned C-Tree state map and let \(d(v)\) be node depth. Suppress the C1/C2/C3
channel index for notation; the same construction applies to a fixed weighted
mixture of law channels. Let \(\gamma\ge 0\) be the declared depth-discount
weight. The oracle and proxy node losses are
\[
  \ell_v^\ast=\ell_{\mathrm{law}}(v;f^\ast,g),
  \qquad
  \widehat\ell_v=\ell_{\mathrm{law}}(v;\widehat f,g).
\]
Here \(\ell_{\mathrm{law}}\) is the per-node local-law discrepancy,
\(f^\ast\) is the target oracle readout, and \(\widehat f\) is the fitted
proxy readout. The starred loss is oracle-valued and directly observed only
when the node is sampled for an oracle call; the hatted loss is the proxy
estimate observed densely.
The oracle local-law target is
\begin{equation}
\label{eq:v8-app-oracle-law-target}
  L_{\mathrm{law}}^\ast(g,T)
  =
  \sum_{v\in V(T)}
  \gamma^{d(v)}\ell_v^\ast .
\end{equation}
This is the oracle-relative local-law channel for the same learned \(g\) that
would appear in the objective if every node could be evaluated with
\(f^\ast\).

\subsection{Node-Level Correction}
\label{app:v8-node-correction}

Let \(R_v\in\{0,1\}\) be the node-oracle observation indicator:
\[
  R_v=
  \begin{cases}
    1, & \text{if node }v\text{ is sampled and evaluated with }f^\ast,\\
    0, & \text{otherwise.}
  \end{cases}
\]
Let \(\pi_v>0\) be the logged marginal probability that node \(v\) would be
sampled under the audit design. The ratio \(R_v/\pi_v\) is therefore the
inverse-propensity weight attached to an observed oracle residual. The
corrected node loss is
\begin{equation}
\label{eq:v8-app-node-correction}
  \ell_v^{\mathrm{corr}}
  =
  \widehat\ell_v
  +
  \frac{R_v}{\pi_v}
  \left(\ell_v^\ast-\widehat\ell_v\right).
\end{equation}
Equivalently, with \(b_v=\widehat\ell_v-\ell_v^\ast\),
\begin{equation}
\label{eq:v8-app-node-bias-form}
  \ell_v^{\mathrm{corr}}
  =
  \ell_v^\ast
  +
  \left(1-\frac{R_v}{\pi_v}\right)b_v .
\end{equation}
The correction keeps the residual bias explicit. If the logged propensity
matches the true marginal inclusion probability, the conditional expectation of
the residual term is zero. If the proxy local-law loss is already exact
\((b_v=0)\), the residual is zero for any observation state.

\subsection{Tree Aggregation}
\label{app:v8-tree-correction}

The corrected local-law tree loss is
\begin{equation}
\label{eq:v8-app-corrected-tree-loss}
  L_{\mathrm{law}}^{\mathrm{corr}}(g,T)
  =
  \sum_{v\in V(T)}
  \gamma^{d(v)}
  \left[
    \widehat\ell_v
    +
    \frac{R_v}{\pi_v}
    \left(\ell_v^\ast-\widehat\ell_v\right)
  \right].
\end{equation}
Substituting the bias form gives
\[
  L_{\mathrm{law}}^{\mathrm{corr}}(g,T)
  =
  L_{\mathrm{law}}^\ast(g,T)
  +
  \sum_{v\in V(T)}
  \gamma^{d(v)}
  \left(1-\frac{R_v}{\pi_v}\right)b_v .
\]
Under correct logged propensities, this aggregate is conditionally unbiased for
the oracle law target. Under misspecified propensities, the remaining bias is
the discounted sum of the displayed residuals.

\subsection{Design-Based Coverage}
\label{app:v8-design-coverage}

The IPW/HT role is estimation, not target definition. For a finite population
\(\mathcal U\), let \(Y_i\) be the outcome or loss contribution for unit
\(i\), let \(Z_i\in\{0,1\}\) indicate whether unit \(i\) is sampled, and let
\(\pi_i>0\) be its logged inclusion probability. The
inverse-propensity-weighted (IPW) Horvitz--Thompson (HT) mean is
\[
  \widehat\mu_{\mathrm{HT}}
  =
  \frac{1}{|\mathcal U|}
  \sum_{i\in\mathcal U}\frac{Z_i}{\pi_i}Y_i .
\]
When the logged propensities match the sampling design, this estimates the
corresponding full-population component. For root supervision,
coverage-normalized objectives divide by the selected root-label count or use
the equivalent constant-inclusion HT mean. For local laws, the sampled
\(f^\ast\) node calls identify the oracle residuals in
Equation~\eqref{eq:v8-app-node-correction}, while \(\widehat f\) supplies the
dense proxy losses.

Coverage therefore changes sampling variance and effective sample size, not
the population root/local tradeoff. Here effective sample size is the
variance-facing size of the weighted sample after uneven propensities or
clipping.

\subsection{Objective Interface}
\label{app:v8-objective-interface}

The ideal objective has one root channel and one oracle local-law channel:
\[
  J_\lambda^\ast(g;T)
  =
  (1-\lambda)L_{\mathrm{root}}(g)
  +
  \lambda L_{\mathrm{law}}^\ast(g,T).
\]
In finite data, the corrected plug-in objective has one root channel and one
corrected local-law channel:
\begin{equation}
\label{eq:v8-app-final-objective}
  J_\lambda^{\mathrm{corr}}(g;T)
  =
  (1-\lambda)L_{\mathrm{root}}(g)
  +
  \lambda L_{\mathrm{law}}^{\mathrm{corr}}(g,T).
\end{equation}
In lambda mode, the analyst supplies \(\lambda\in[0,1]\) and an active law
set. The root channel receives weight \(1-\lambda\), and the active local laws
split \(\lambda\) by a fixed predeclared rule.

In explicit-weight mode, the analyst instead supplies nonnegative weights
\[
  w_{\mathrm{root}},w_1,\ldots,w_k
\]
for the root channel and the \(k\) active local-law channels. The same
corrected objective is written as
\[
  J(g;T)
  =
  \frac{w_{\mathrm{root}}}{W}L_{\mathrm{root}}(g)
  +
  \sum_{j=1}^k
  \frac{w_j}{W}L_j^{\mathrm{corr}}(g,T),
  \qquad
  W=w_{\mathrm{root}}+\sum_{j=1}^k w_j .
\]
Here \(w_{\mathrm{root}}\) is the root-channel weight,
\(L_{\mathrm{root}}\) is the root loss, \(w_j\) is the weight for corrected
local-law channel \(j\), \(L_j^{\mathrm{corr}}\) is that channel's corrected
tree loss, \(k\) is the number of active local-law channels, and \(W\) is the
total normalizing weight. The implied local-law share is \(\sum_j w_j/W\). It
is reported for interpretation and is not supplied as an additional input.
Lambda mode and explicit-weight mode are alternative interfaces; a single
mathematical objective should not take both \(\lambda\) and explicit channel
weights as independent controls.

The proxy-oracle gap, propensity residual, standard error, and clipping
envelope describe the quality of \(L_{\mathrm{law}}^{\mathrm{corr}}\). They are
reported with the audit certificate and do not enter
Equation~\eqref{eq:v8-app-final-objective} as separate optimized loss terms.
Here the propensity residual is the displayed remaining bias term, the standard
error is the sampling uncertainty of the corrected aggregate, and the clipping
envelope is any declared bound applied to weights or residuals.

\subsection{Certificate Decomposition}
\label{app:v8-certificate-decomposition}

For a sampled node-level certificate, let \(\mathcal U\) be the finite
population of realized tree units, \(N=|\mathcal U|\), \(Z_i\) the sampled-unit
indicator, and \(\pi_i>0\) the logged marginal inclusion probability. Let
\[
  D_i^J=d_Y(J(s_i),J(S(i)))
\]
be the accessible-judge distortion for state \(s_i\) against the represented
span \(S(i)\). The unclipped HT distortion estimate is
\[
  \widehat{\Delta}_{\mathrm{HT}}^J
  =
  \frac{1}{N}
  \sum_{i\in\mathcal U}
  \frac{Z_i}{\pi_i}D_i^J .
\]
The corresponding full-population accessible-judge mean is
\[
  \Delta_J
  =
  \frac{1}{N}\sum_{i\in\mathcal U}D_i^J .
\]
Correct logged propensities give
\(\mathbb E[\widehat{\Delta}_{\mathrm{HT}}^J\mid\mathcal U]=\Delta_J\).

The reported certificate uses four terms:
\[
  |G_{\mathrm{meth}}|
  \le
  C_{\mathrm{meth}}\widehat{\Delta}_{\mathrm{HT}}
  +
  B_{\mathrm{cal}}
  +
  B_{\mathrm{est}}
  +
  B_{\mathrm{clip}} .
\]
Here \(C_{\mathrm{meth}}\) transports oracle distortion to the downstream
method gap, \(B_{\mathrm{cal}}\) bounds the accessible-judge/oracle mismatch,
\(B_{\mathrm{est}}\) bounds sampling error around the HT mean after method
transport, and \(B_{\mathrm{clip}}\) bounds bias from any declared clipping
rule. In a high-probability report, each envelope carries its own event
probability and the total failure probability is the declared union-bound
sum.
